% Faithful typesetting transcription of the formal CCC2022.23 publication, Sections2 and AppendixA.
\section*{Author prerequisites (Section2)}
\paragraph{Notation.} We define $[n]\triangleq\{0,1,\ldots,n-1\}$, $B_{n,m}$ as the set of all functions from $\mathbb F_2^n$ to $\mathbb F_2^m$ and $B_n\triangleq B_{n,1}$. The Hamming weight of $x\in\mathbb F_2^n$, denoted by $|x|$, is defined as the number of 1's in $x$; the Hamming distance $\Delta(x,y)$ of $x$ and $y$ from $\mathbb F_2^n$ is defined as the Hamming weight of the bitwise XOR of them. The concatenation of two strings $x$ and $y$ is denoted by $x\Vert y$, and the $i$th bit of $x$ (0-indexed) is denoted by $x_i$. Graphs are undirected by default. We assume that all functions used as parameters of our constructions are poly-logarithmic time constructible, i.e., for any function $\ell(n)$ that is used as a parameter in our results, there is a polynomial-time algorithm $A$ such that $A(n)$ outputs the binary representation of $\ell(n)$.
We use $x\leftarrow D$ to denote a random element $x$ sampled from a distribution $D$, and $D(x_0)\triangleq\Pr_{x\leftarrow D}[x=x_0]$. Natural numbers are represented in binary when being fed into Turing machines or circuits; and $1^n$ represents the unary representation of $n$. We assume basic familiarity with cryptographic primitives like one-way functions and pseudorandom functions, and typical complexity classes like $\mathsf P$, $\mathsf{NP}$, and $\oplus\mathsf P$ (see, e.g., [3]).
\subsection*{2.1 Probability Theory}
Let $U(S)$ be the uniform distribution supported on a set $S$, and $U_\ell\triangleq U(\{0,1\}^\ell)$. A family $D=\{D_n\}_{n\ge1}$ of distributions is said to be samplable with seed length $\ell(n)$ (or $\ell(n)$-samplable) if there is a polynomial-time algorithm $A$ such that $A(1^n,U_{\ell(n)})$ samples $D_n$. For every $s\in\{0,1\}^\ell$, we say $A(1^n,s)$ is the element corresponding to the seed $s$.
The statistical distance between two distributions $D_1$ and $D_2$ over a set $S$, denoted by $\operatorname{SD}(D_1,D_2)$, is defined as
\[\operatorname{SD}(D_1,D_2)\triangleq\frac12\sum_{u\in S}|D_1(u)-D_2(u)|.\]
A distribution $D_2$ is said to be $\delta$-close to $D_1$ if their statistical distance is at most $\delta$. Statistical distance characterizes the intractability of distinguishing two distributions by any test (even computationally unbounded): if $D_1$ and $D_2$ have statistical distance at most $\delta$, then for any stochastic process $T$, $|\Pr_{u\leftarrow D_1,T}[T(u)=1]-\Pr_{u\leftarrow D_2,T}[T(u)=1]|\le\delta$.
\subsection*{2.2 Circuit Classes}
In this paper we will work with various circuit classes. In general, a circuit is an acyclic graph where each of its nodes can be an input variable, a constant $c\in\{0,1\}$, or a gate. One or more nodes are marked as output nodes (together with an index $i$ denoting the corresponding output bit). The depth of a circuit is the number of edges on the longest path from any input variable to an output node. An $n$-input $m$-output circuit computes a function in $B_{n,m}$.
\begin{description}
\item[{$B_2$ circuits.}] A $B_2$ circuit (or general circuit) contains fan-in-2 gates that can compute any binary function $f\in B_2$. The size of a $B_2$ circuit refers to the number of gates involved.
\item[{$\mathsf{NC}^1$ circuits.}] An $\mathsf{NC}^1$ circuit is a $B_2$ circuit with $O(\log n)$ depth. Note that a single-output $\mathsf{NC}^1$ circuit of depth $d$ can be converted into a formula of size $O(2^d)$.
\item[{$\mathsf{CC}^0[2]$ circuits.}] A $\mathsf{CC}^0[2]$ circuit is a constant-depth circuit with only XOR gates of unbounded fan-in. It is easy to see that $\mathsf{CC}^0[2]$ circuits can only compute linear functions over $\mathbb F_2$. The complexity of a $\mathsf{CC}^0[2]$ circuit is usually measured by the number of wires.
\item[{$\mathsf{AC}^0[m]$ circuits.}] An $\mathsf{AC}^0[m]$ circuit is a constant-depth circuit with fan-in-1 NOT gates and unbounded fan-in gates over $\{\mathrm{AND},\mathrm{OR},\mathrm{MOD}_m\}$, where $\mathrm{MOD}_m(x_1,\ldots,x_k)=1$ if and only if $x_1+x_2+\cdots+x_k\equiv0\pmod m$. Similar to $\mathsf{CC}^0[2]$ circuits, the complexity of an $\mathsf{AC}^0[m]$ circuit is measured by the number of wires.
\end{description}
A probabilistic $B_2$ circuit (or simply a probabilistic circuit) with input size $n$ and output size $m$ is a distribution $C_n$ over $n$-input $m$-output $B_2$ circuits. The circuit complexity of a probabilistic circuit is defined as the maximum complexity of functions in its support. A family $C=\{C_n\}_{n\ge1}$ of $n$-input 1-output probabilistic circuit is said to decide a language $L$ with error probability $\varepsilon=\varepsilon(n)$ if for sufficiently large $n$ and all $x\in\mathbb F_2^n$, $\Pr_{C\leftarrow C_n}[C(x)\ne L(x)]\le\varepsilon$.
\subsection*{2.3 Hash and 1-detector}
\begin{authorclaim}{Definition2.1 (Almost universal hash function)}
Let $m=m(n)$ and $\varepsilon=\varepsilon(n)$ be two parameters. An $m$-output $\varepsilon$-almost universal hash function is a family of distributions $H=\{H_n\}_{n\ge1}$, where $H_n$ is supported on $B_{n,m}$, such that $\Pr_{h\leftarrow H_n}[h(x)=h(y)]\le\varepsilon(n)$ for all $x\ne y$ and sufficiently large $n$. The parameter $\varepsilon$ is called its collision probability. It is linear if every function in the support of $H_n$ is linear. It is said to be $s(n)$-samplable if the family $H$ of distributions is samplable with seed length $s(n)$.
\end{authorclaim}
We will make heavy use of the notion of 1-detectors in [13], as parts of our almost universal hash construction.
\begin{authorclaim}{Definition2.2 (1-detector)}
An $m$-output (randomized) 1-detector with range $r$ and error $\varepsilon$ is a family of distributions $D=\{D_n\}_{n\ge1}$, where $D_n$ is supported on linear functions in $B_{n,m}$, such that for sufficiently large $n$, $\Pr_{L\leftarrow D_n}[L(x)\ne0]\le\varepsilon(n)$ for all $x\in\mathbb F_2^n$ satisfying $0<|x|\le r$. It is said to be $s(n)$-samplable if the family $D$ of distributions is samplable with seed length $s(n)$.
\end{authorclaim}
We emphasize that the definition requires the functions in the family of distributions to be linear. This ensure that, for any randomized 1-detector $D_n$ with range $r$ and error $\varepsilon$, and any $x_1\ne x_2$ with Hamming distance at most $r$, it holds that $\Pr_{L\leftarrow D_n}[L(x_1)\ne L(x_2)]\le\varepsilon(n)$. The shrinkage of a hash function (or 1-detector) is defined as the input length $n$ divided by the output length $m$ (e.g., poly-logarithmic shrinkage means $m=n/\operatorname{polylog}(n)$).
\subsection*{2.4 $\varepsilon$-biased set and $k$-wise independence}
To derandomize the low-complexity hash function in [13], we need several standard derandomization tools including strongly explicit $\varepsilon$-biased sets, $k$-wise independent distributions, and $k$-wise independent permutations.
\begin{authorclaim}{Definition2.3 ($\varepsilon$-biased set)}
For any $n\ge1$ and $\varepsilon\in(0,1/2)$, a set $S\subseteq\mathbb F_2^n$ is said to be $\varepsilon$-biased if for all non-zero $v\in\mathbb F_2^n$, $\Pr_{w\leftarrow S}[\langle w,v\rangle=0]\in[1/2-\varepsilon,1/2+\varepsilon]$.
\end{authorclaim}
\begin{authorclaim}{Theorem2.4 (Alon, Goldreich, H\aa stad, and Peralta [1])}
For any constant $\varepsilon\in(0,1/2)$, there is a family $\{S_n\subseteq\mathbb F_2^n\}_{n\ge1}$ of $\varepsilon$-biased sets such that $|S_n|=\widetilde O(n^2)$. Moreover, there is an algorithm $A(n,i,j)$ running in time $\operatorname{poly}(\log n)$ that computes the $j$th bit of the $i$th vector in $S_n$.
\end{authorclaim}
\begin{authorclaim}{Definition2.5}
A $k$-wise $\varepsilon$-dependent distribution with length $n$ and alphabet size $p$ is a distribution $D$ over $[p]^n$ such that for all $0\le i_1<i_2<\cdots<i_k<n$, the distribution over $[p]^k$ obtained by restricting $D$ to the $i_1$th, $i_2$th, $\ldots$, $i_k$th coordinates is $\varepsilon$-close to the uniformly distribution over $[p]^k$. It is said to be a $k$-wise independent distribution if $\varepsilon=0$.
\end{authorclaim}
\begin{authorclaim}{Theorem2.6 (Alon, Goldreich, H\aa stad, and Peralta [1])}
Let $\varepsilon=\varepsilon(n)>0$ be a parameter. There is an algorithm $A$ such that $A(1^n,k,U_r)$ runs in time $\operatorname{poly}(n,\log(1/\varepsilon))$ and samples a $k$-wise $\varepsilon$-dependent distribution with length $n$ and alphabet size 2 for every $1\le k\le n$, where $r=O(k+\log\log n+\log(1/\varepsilon))$. Moreover, there is an algorithm $B$ such that $B(n,k,v,i)$ runs in $\operatorname{poly}(\log n,k,\log(1/\varepsilon))$ time and computes the $i$th coordinate of the string corresponding to the seed $v$.
\end{authorclaim}
\begin{authorclaim}{Definition2.7}
Let $n\ge1$ be the number of elements, and $P_n$ be the set of all permutations over $[n]$. A $k$-wise $\varepsilon$-dependent permutation is a distribution $D$ over $P_n$ such that for all $0\le i_1<i_2<\cdots<i_k<n$, the distribution of $(f(i_1),f(i_2),\ldots,f(i_k))$ with $f\leftarrow D$ is $\varepsilon$-close to the uniform distribution over $\{(j_1,j_2,\ldots,j_k)\mid j_1,\ldots,j_k\in[n]\text{ are pairwisely distinct}\}$.
\end{authorclaim}
\begin{authorclaim}{Theorem2.8 (Kaplan, Naor, and Reingold [25], Theorem5.9)}
Let $n$ be a power of 2, $2\le k\le n$, and $\varepsilon>2^{-n}$. There is an $n$-element $k$-wise $\varepsilon$-dependent permutation that is samplable with seed length $O(k\log n+\log(1/\varepsilon))$. Moreover, there is an algorithm $A(n,v,i)$ that runs in $\operatorname{poly}(\log n,k,\log(1/\varepsilon))$ time and outputs $\rho_v(i)$ for the permutation $\rho_v$ corresponding to the seed $v$.
\end{authorclaim}
\subsection*{2.5 Expander Graphs}
We will need strongly explicit expander graphs in our proofs. We first recall the definition of expander graphs.
\begin{authorclaim}{Definition2.9}
An $n$-vertex $d$-regular graph $G$ is called an $(n,d,\lambda)$ expander graph (or $(n,d,\lambda)$-graph) if $\lambda_2(G)\le\lambda$, where $\lambda_2(G)$ denotes the second largest eigenvalue of normalized adjacency matrix (i.e., adjacency matrix divided by $d$) of $G$. A family of graphs $\{G_n\}_{n\ge1}$ is an expander graph family if there exists constants $d$ and $\lambda<1$ such that for all sufficiently large $n$, $G_n$ is an $(n,d,\lambda)$-graph.
\end{authorclaim}
We will make use of the following construction of strongly explicit expander graphs.
\begin{authorclaim}{Theorem2.10 (Strongly explicit expander; see [3, Theorem21.19])}
There exists an expander graph family $\{G_n\}_{n\ge1}$ and an algorithm $A(n,v,i)$ that runs in $\operatorname{polylog}(n)$ time (i.e., polynomial in input length) and outputs the $i$th neighbor of the node $v$ in $G_n$, where $v\in[n]$ and $i\in[d]$.
\end{authorclaim}
Performing random walk on expanders is a standard derandomization technique. Consider the task to find a good element from $n$ elements in which there is a constant fraction of them being good. A trivial approach is to sample $\ell$ independently random elements, which has $\exp(-\Omega(\ell))$ error probability but requires $\ell\log n$ random bits. By applying the following lemma, we can reduce the randomness complexity to $O(\log n+\ell)$ while keeping the error probability to be exponentially small.
\begin{authorclaim}{Lemma2.11 (Expander walk; see [3, Theorem21.12])}
Let $G$ be an $(n,d,\lambda)$-graph and $S$ be a subset of vertices of size at most $\beta n$ for some $\beta\in(0,1)$. Let $X_1,X_2,\ldots,X_k$ be random variables denoting a random walk in $G$ (where $X_1$ is uniformly chosen), then
\[\Pr[\forall1\le i\le k, X_i\in S]\le((1-\lambda)\sqrt\beta+\lambda)^{k-1}.\]
\end{authorclaim}
\subsection*{2.6 Graph with large girth}
The girth of an undirected graph is the length of the minimum cycle in it. We need the following construction of strongly explicit graphs with large girth.
\begin{authorclaim}{Theorem2.12 (Adapted from [27]; see AppendixA)}
Let $r=r(n)=n^{o(1)}$ be a parameter. For every sufficiently large $n$, there exists an $m=\Theta(n/r)$ and a regular graph $G_{m,n}$ with $m$ vertices, $n$ edges, and girth $\Omega(\log n/\log r)$. Moreover, there exists a $\operatorname{polylog}(n)$-time algorithm $A(n,i)$ for $i\in[n]$ that outputs the indices of the two endpoints of the $i$th edge in $G_{m,n}$, and a $\operatorname{polylog}(n)$-time algorithm $B(n,i,j)$ for $i\in[m]$ that outputs the $j$th edge attaching to the $i$th vertex.
\end{authorclaim}
